% ═══════════════════════════════════════════════════════════════
% SLIDES: HAY (Repaso) + ESTAR
% Klasse 7 | Unidad 2 | Mi instituto | DS 02
% ═══════════════════════════════════════════════════════════════

\documentclass[aspectratio=169, 11pt]{beamer}

\usepackage[utf8]{inputenc}
\usepackage[T1]{fontenc}
\usepackage[ngerman]{babel}
\usepackage{palatino}
\usepackage{helvet}
\usepackage{tabularx}
\usepackage{booktabs}
\usepackage{array}
\usepackage{tikz}
\usepackage{amssymb}

% ═══════════════════════════════════════════════════════════════
% THEME & FARBEN
% ═══════════════════════════════════════════════════════════════

\usetheme{default}
\usecolortheme{default}

% Spanisch-Farben
\definecolor{spanischrot}{HTML}{C41E3A}
\definecolor{dunkelgrau}{HTML}{333333}
\definecolor{hellgrau}{HTML}{F5F5F5}
\definecolor{erfolg}{HTML}{2E7D32}
\definecolor{fehler}{HTML}{C62828}
\definecolor{blau}{HTML}{1565C0}

% Beamer-Anpassungen
\setbeamercolor{frametitle}{fg=white, bg=spanischrot}
\setbeamercolor{title}{fg=white, bg=spanischrot}
\setbeamercolor{structure}{fg=spanischrot}
\setbeamercolor{block title}{fg=white, bg=spanischrot}
\setbeamercolor{block body}{fg=dunkelgrau, bg=hellgrau}
\setbeamercolor{item}{fg=spanischrot}
\setbeamercolor{subitem}{fg=spanischrot}
\setbeamercolor{alerted text}{fg=spanischrot}

\setbeamerfont{frametitle}{series=\bfseries, size=\Large}
\setbeamerfont{title}{series=\bfseries, size=\huge}

% Keine Navigation
\setbeamertemplate{navigation symbols}{}
\setbeamertemplate{footline}[frame number]

% ═══════════════════════════════════════════════════════════════
% BEFEHLE
% ═══════════════════════════════════════════════════════════════

% Spanisch hervorheben
\newcommand{\es}[1]{\textbf{\textcolor{spanischrot}{#1}}}
% Deutsch (klein, grau)
\newcommand{\de}[1]{{\small\textcolor{dunkelgrau}{\textit{#1}}}}
% Richtig/Falsch
\newcommand{\richtig}[1]{\textcolor{erfolg}{\checkmark\,#1}}
\newcommand{\falsch}[1]{\textcolor{fehler}{$\times$\,#1}}
% Blau für ESTAR
\newcommand{\estar}[1]{\textbf{\textcolor{blau}{#1}}}

% ═══════════════════════════════════════════════════════════════
% DOKUMENT
% ═══════════════════════════════════════════════════════════════

\begin{document}

% ---------------------------------------------------------------
% TITELFOLIE
% ---------------------------------------------------------------
{
\setbeamercolor{background canvas}{bg=spanischrot}
\begin{frame}[plain]
\begin{center}
\vspace{2cm}
{\Huge\bfseries\textcolor{white}{Hay y estar}}\\[0.5cm]
{\Large\textcolor{white}{Wo ist was? Was gibt es?}}\\[1.5cm]
{\large\textcolor{white}{Spanisch · Klasse 7 · Unidad 2}}
\end{center}
\end{frame}
}

% ---------------------------------------------------------------
% LERNZIELE
% ---------------------------------------------------------------
\begin{frame}{Heute lernst du...}
\begin{itemize}
    \item[\textcolor{spanischrot}{\textbullet}] \textbf{Repaso:} \es{hay} -- alle Verwendungen
    \item[\textcolor{spanischrot}{\textbullet}] \textbf{Neu:} \estar{estar} = sein (Ort/Zustand)
    \item[\textcolor{spanischrot}{\textbullet}] \textbf{Neu:} Präpositionen (encima de, debajo de, ...)
    \item[\textcolor{spanischrot}{\textbullet}] \textbf{Unterschied:} \es{hay} vs \estar{estar}
\end{itemize}

\vspace{0.8cm}

\begin{block}{Am Ende kannst du...}
...sagen, WAS ES GIBT und WO ETWAS IST!
\end{block}
\end{frame}

% ---------------------------------------------------------------
% WARM-UP
% ---------------------------------------------------------------
\begin{frame}{Warm-up: ¿Qué hay en tu mochila?}

\begin{center}
{\Large Erzähle deinem Partner:}\\[0.5cm]
{\Huge \es{En mi mochila hay...}}
\end{center}

\vspace{1cm}

\begin{block}{Redemittel}
\es{En mi mochila hay un libro, dos cuadernos, muchos bolígrafos y libros.}\\
{\small\textcolor{dunkelgrau}{(un libro = ein Buch, dos cuadernos = zwei Hefte, muchos = viele, libros = Bücher)}}
\end{block}

\end{frame}

% ---------------------------------------------------------------
% HAY REPASO: ALLE 4 FÄLLE
% ---------------------------------------------------------------
\begin{frame}{\es{hay} -- Wiederholung (4 Fälle)}

\begin{center}
{\Huge \es{hay} = es gibt}
\end{center}

\vspace{0.5cm}

\begin{columns}[T]
\begin{column}{0.48\textwidth}
\textbf{1. + un/una}\\
\es{Hay \textbf{una} mesa.}\\[0.5em]

\textbf{2. + Zahl}\\
\es{Hay \textbf{tres} sillas.}\\
\end{column}

\begin{column}{0.48\textwidth}
\textbf{3. + mucho/poco}\\
\es{Hay \textbf{muchos} libros.}\\
\es{Hay \textbf{pocos} cuadernos.}\\[0.5em]

\textbf{4. + ohne Artikel}\\
\es{Hay \textbf{libros}.}\\
\es{¿Hay \textbf{leche}?}\\
\end{column}
\end{columns}

\vspace{0.5cm}

\begin{alertblock}{Merke}
Nach \es{hay} kommt \textbf{NIE} el/la!
\end{alertblock}

\end{frame}

% ---------------------------------------------------------------
% HAY: ÜBUNG AB 1-2
% ---------------------------------------------------------------
\begin{frame}{$\rightarrow$ Arbeitsblatt: Aufgabe 1 + 2}

\begin{center}
{\Large Bearbeite jetzt \textbf{Aufgabe 1} und \textbf{Aufgabe 2}!}
\end{center}

\vspace{0.8cm}

\begin{block}{Aufgabe 1}
Ordne die HAY-Regeln den Beispielen zu.
\end{block}

\begin{block}{Aufgabe 2}
Bilde Sätze mit \es{hay} (verschiedene Strukturen).
\end{block}

\vspace{0.5cm}

\begin{center}
{\small\textcolor{dunkelgrau}{Zeit: ca. 8 Minuten}}
\end{center}

\end{frame}

% ---------------------------------------------------------------
% ESTAR: EINFÜHRUNG
% ---------------------------------------------------------------
\begin{frame}{\estar{estar} = sein (Ort/Zustand)}

\begin{center}
{\Huge \estar{estar}}\\[0.3cm]
{\Large = sein \de{(für Ort/Zustand)}}
\end{center}

\vspace{0.5cm}

\begin{columns}[c]
\begin{column}{0.45\textwidth}
\begin{center}
\begin{tabular}{@{}ll@{}}
yo & \estar{estoy} \\[0.3em]
tú & \estar{estás} \\[0.3em]
él/ella & \estar{está} \\
\end{tabular}
\end{center}
\end{column}

\begin{column}{0.45\textwidth}
\begin{center}
\begin{tabular}{@{}ll@{}}
nosotros & \estar{estamos} \\[0.3em]
vosotros & \estar{estáis} \\[0.3em]
ellos & \estar{están} \\
\end{tabular}
\end{center}
\end{column}
\end{columns}

\vspace{0.5cm}

\begin{alertblock}{Merke}
Alle Formen außer \estar{estamos} haben Akzent!
\end{alertblock}

\end{frame}

% ---------------------------------------------------------------
% PRÄPOSITIONEN
% ---------------------------------------------------------------
\begin{frame}{Preposiciones \de{-- Präpositionen}}

\begin{columns}[T]
\begin{column}{0.48\textwidth}
\begin{tabular}{@{}ll@{}}
\estar{en} & \de{in, auf, an} \\[0.6em]
\estar{encima de} & \de{auf, oberhalb} \\[0.6em]
\estar{debajo de} & \de{unter} \\[0.6em]
\estar{al lado de} & \de{neben} \\
\end{tabular}
\end{column}

\begin{column}{0.48\textwidth}
\begin{tabular}{@{}ll@{}}
\estar{delante de} & \de{vor} \\[0.6em]
\estar{detrás de} & \de{hinter} \\[0.6em]
\estar{entre} & \de{zwischen} \\
\end{tabular}
\end{column}
\end{columns}

\vspace{0.8cm}

\begin{block}{Beispiele}
\estar{El libro está \textbf{encima de} la mesa.}\\
\estar{La mochila está \textbf{debajo de} la silla.}
\end{block}

\end{frame}

% ---------------------------------------------------------------
% ESTAR: ÜBUNG AB 3
% ---------------------------------------------------------------
\begin{frame}{$\rightarrow$ Arbeitsblatt: Aufgabe 3}

\begin{center}
{\Large Bearbeite jetzt \textbf{Aufgabe 3}!}
\end{center}

\vspace{0.8cm}

\begin{block}{Aufgabe 3}
Ergänze \estar{estar} und die passende Präposition.
\end{block}

\vspace{0.5cm}

\begin{center}
{\large Frage: \estar{¿Dónde está...?} \de{Wo ist...?}}
\end{center}

\vspace{0.5cm}

\begin{center}
{\small\textcolor{dunkelgrau}{Zeit: ca. 8 Minuten}}
\end{center}

\end{frame}

% ---------------------------------------------------------------
% HAY VS ESTAR: KONTRAST
% ---------------------------------------------------------------
\begin{frame}{\es{hay} vs \estar{estar} -- Der Unterschied}

\begin{columns}[T]
\begin{column}{0.48\textwidth}
\begin{center}
{\Large\bfseries\textcolor{spanischrot}{hay}}\\[0.3cm]
{\large Existenz}\\[0.5em]
\de{Gibt es das?}\\[0.5em]
(unbekannt/unbestimmt)
\end{center}

\vspace{0.5cm}

\es{Hay \textbf{un} libro en la mesa.}\\[0.3em]
\de{Es gibt \textbf{ein} Buch.}\\[0.3em]
{\small (irgendein Buch)}

\vspace{0.5cm}

$\rightarrow$ \textbf{un/una}, Zahlen, muchos
\end{column}

\begin{column}{0.48\textwidth}
\begin{center}
{\Large\bfseries\textcolor{blau}{estar}}\\[0.3cm]
{\large Ort}\\[0.5em]
\de{Wo ist es?}\\[0.5em]
(bekannt/bestimmt)
\end{center}

\vspace{0.5cm}

\estar{\textbf{El} libro está en la mesa.}\\[0.3em]
\de{\textbf{Das} Buch ist auf dem Tisch.}\\[0.3em]
{\small (ein bestimmtes Buch)}

\vspace{0.5cm}

$\rightarrow$ \textbf{el/la}, mi/tu/su
\end{column}
\end{columns}

\end{frame}

% ---------------------------------------------------------------
% HAY VS ESTAR: BEISPIELE
% ---------------------------------------------------------------
\begin{frame}{Beispiele: \es{hay} vs \estar{estar}}

\begin{center}
\begin{tabular}{@{}p{6.5cm}p{6.5cm}@{}}
\es{¿Hay una goma?} & \de{Gibt es einen Radierer?} \\[0.5em]
\es{Sí, hay una goma.} & \de{Ja, es gibt einen.} \\[1em]
\estar{¿Dónde está la goma?} & \de{Wo ist der Radierer?} \\[0.5em]
\estar{La goma está en el estuche.} & \de{Der Radierer ist in der Federtasche.} \\
\end{tabular}
\end{center}

\vspace{0.8cm}

\begin{block}{Eselsbrücke}
\textbf{hay} reimt auf \textit{dabei} → „Was ist alles dabei?"\\
\textbf{ESTAR} = \textbf{E}xakter \textbf{STA}ndo\textbf{R}t \de{(Wo genau?)}
\end{block}

\end{frame}

% ---------------------------------------------------------------
% HAY VS ESTAR: ÜBUNG AB 4
% ---------------------------------------------------------------
\begin{frame}{$\rightarrow$ Arbeitsblatt: Aufgabe 4}

\begin{center}
{\Large Bearbeite jetzt \textbf{Aufgabe 4}!}
\end{center}

\vspace{0.8cm}

\begin{block}{Aufgabe 4}
Entscheide: \es{hay} oder \estar{estar}?\\
Kreuze an und ergänze die richtige Form.
\end{block}

\vspace{0.5cm}

\begin{center}
{\large Denke nach: \textbf{Existenz} oder \textbf{Ort}?}
\end{center}

\vspace{0.5cm}

\begin{center}
{\small\textcolor{dunkelgrau}{Zeit: ca. 8 Minuten}}
\end{center}

\end{frame}

% ---------------------------------------------------------------
% ANWENDUNG: AB 5
% ---------------------------------------------------------------
\begin{frame}{$\rightarrow$ Arbeitsblatt: Aufgabe 5}

\begin{center}
{\Large Beschreibe deinen \textbf{Klassenraum}!}
\end{center}

\vspace{0.5cm}

\begin{columns}[T]
\begin{column}{0.48\textwidth}
\begin{block}{4 Sätze mit \es{hay}}
\es{En mi clase hay...}\\
\es{Hay muchos/as...}\\
\es{Hay un/una...}
\end{block}
\end{column}

\begin{column}{0.48\textwidth}
\begin{block}{4 Sätze mit \estar{estar}}
\estar{El/La ... está...}\\
\estar{Mi ... está...}\\
\estar{Los/Las ... están...}
\end{block}
\end{column}
\end{columns}

\vspace{0.5cm}

\begin{center}
{\small\textcolor{dunkelgrau}{Zeit: ca. 8 Minuten}}
\end{center}

\end{frame}

% ---------------------------------------------------------------
% SPIEL: EINFÜHRUNG
% ---------------------------------------------------------------
\begin{frame}
\begin{tikzpicture}[remember picture, overlay]
    \fill[blau] (current page.north west) rectangle (current page.south east);
\end{tikzpicture}
\begin{center}
\vspace{1.5cm}
{\Huge\bfseries\textcolor{white}{¡Vamos a jugar!}}\\[1cm]
{\Large\textcolor{white}{Adivina dónde está}}\\[0.5cm]
{\large\textcolor{white}{Rate, wo es ist!}}
\end{center}
\end{frame}

% ---------------------------------------------------------------
% SPIEL: REGELN
% ---------------------------------------------------------------
\begin{frame}{Spielregeln: Adivina dónde está}

\begin{block}{So geht's:}
\begin{enumerate}
    \item Ein Schüler schließt die Augen.\\
    {\footnotesize\textcolor{dunkelgrau}{\textit{Un alumno cierra los ojos.}}}
    \item Die Klasse versteckt einen Gegenstand.\\
    {\footnotesize\textcolor{dunkelgrau}{\textit{La clase esconde un objeto.}}}
    \item Der Schüler fragt: \estar{¿Dónde está el/la...?}
    \item Die Klasse gibt Hinweise mit Präpositionen.\\
    {\footnotesize\textcolor{dunkelgrau}{\textit{La clase da pistas con preposiciones.}}}
    \item Nach dem Finden: vollständigen Satz sagen!\\
    {\footnotesize\textcolor{dunkelgrau}{\textit{Después: ¡decir la frase completa!}}}
\end{enumerate}
\end{block}

\end{frame}

% ---------------------------------------------------------------
% SPIEL: REDEMITTEL
% ---------------------------------------------------------------
\begin{frame}{Redemittel für das Spiel}

\begin{columns}[T]
\begin{column}{0.48\textwidth}
\textbf{Fragen:} \de{Preguntas}\\[0.5em]
\estar{¿Dónde está el lápiz?}\\
{\footnotesize\textcolor{dunkelgrau}{\textit{Wo ist der Bleistift?}}}\\[0.4em]
\estar{¿Dónde está la goma?}\\
{\footnotesize\textcolor{dunkelgrau}{\textit{Wo ist der Radierer?}}}\\[0.4em]
\estar{¿Dónde está el estuche?}\\
{\footnotesize\textcolor{dunkelgrau}{\textit{Wo ist die Federtasche?}}}
\end{column}

\begin{column}{0.48\textwidth}
\textbf{Hinweise:} \de{Pistas}\\[0.5em]
\estar{Está \textbf{encima de}...} \de{auf}\\[0.4em]
\estar{Está \textbf{debajo de}...} \de{unter}\\[0.4em]
\estar{Está \textbf{al lado de}...} \de{neben}\\[0.4em]
\estar{Está \textbf{detrás de}...} \de{hinter}\\[0.4em]
\estar{Está \textbf{delante de}...} \de{vor}
\end{column}
\end{columns}

\vspace{0.5cm}

\begin{block}{Nach dem Finden \de{-- Después}}
\estar{¡Sí! El lápiz está debajo de la mesa.}\\
{\footnotesize\textcolor{dunkelgrau}{\textit{Ja! Der Bleistift ist unter dem Tisch.}}}
\end{block}

\end{frame}

% ---------------------------------------------------------------
% SPIEL: RUNDE 1
% ---------------------------------------------------------------
\begin{frame}{Runde 1: ¿Dónde está...?}

\begin{center}
{\Huge\bfseries \estar{¿Dónde está el lápiz?}}\\[0.3em]
{\large\textcolor{dunkelgrau}{\textit{Wo ist der Bleistift?}}}
\end{center}

\vspace{0.8cm}

\begin{block}{Hinweise geben \de{-- Dad pistas}}
\begin{itemize}
    \item \estar{Está encima de...} \de{auf...}
    \item \estar{Está al lado de...} \de{neben...}
    \item \estar{Está detrás de...} \de{hinter...}
\end{itemize}
\end{block}

\end{frame}

% ---------------------------------------------------------------
% SPIEL: RUNDE 2
% ---------------------------------------------------------------
\begin{frame}{Runde 2: ¿Dónde está...?}

\begin{center}
{\Huge\bfseries \estar{¿Dónde está la goma?}}\\[0.3em]
{\large\textcolor{dunkelgrau}{\textit{Wo ist der Radierer?}}}
\end{center}

\vspace{0.8cm}

\begin{block}{Hinweise geben \de{-- Dad pistas}}
\begin{itemize}
    \item \estar{Está debajo de...} \de{unter...}
    \item \estar{Está entre...} \de{zwischen...}
    \item \estar{Está delante de...} \de{vor...}
\end{itemize}
\end{block}

\end{frame}

% ---------------------------------------------------------------
% SPIEL: RUNDE 3
% ---------------------------------------------------------------
\begin{frame}{Runde 3: ¿Dónde está...?}

\begin{center}
{\Huge\bfseries \estar{¿Dónde está el estuche?}}\\[0.3em]
{\large\textcolor{dunkelgrau}{\textit{Wo ist die Federtasche?}}}
\end{center}

\vspace{0.8cm}

\begin{block}{Hinweise geben \de{-- Dad pistas}}
Benutze verschiedene Präpositionen!\\
{\footnotesize\textcolor{dunkelgrau}{\textit{¡Usa diferentes preposiciones!}}}
\end{block}

\end{frame}

% ---------------------------------------------------------------
% SPIEL: VARIANTE
% ---------------------------------------------------------------
\begin{frame}{Variante 2}

\begin{center}
{\Large Lehrer beschreibt einen Ort $\rightarrow$ Schüler erraten den Gegenstand!}\\[0.3em]
{\normalsize\textcolor{dunkelgrau}{\textit{El profesor describe un lugar $\rightarrow$ los alumnos adivinan el objeto.}}}
\end{center}

\vspace{0.6cm}

\begin{block}{Beispiel \de{-- Ejemplo}}
\textbf{Lehrer:} \estar{Está encima de la mesa, al lado del ordenador.}\\
{\footnotesize\textcolor{dunkelgrau}{\textit{Es ist auf dem Tisch, neben dem Computer.}}}\\[0.5em]
\textbf{Schüler:} \es{¡Es el libro!} \de{Das ist das Buch!}
\end{block}

\vspace{0.4cm}

\begin{center}
{\large Wer ist am schnellsten? \textcolor{dunkelgrau}{\textit{¿Quién es el más rápido?}}}
\end{center}

\end{frame}

% ---------------------------------------------------------------
% BONUS: AB AUFGABE 6
% ---------------------------------------------------------------
\begin{frame}{Bonus: Mini-Diálogo}

\begin{center}
{\Large Wer fertig ist: \textbf{Aufgabe Bonus}!}
\end{center}

\vspace{0.5cm}

\begin{block}{Aufgabe}
Schreibe einen Dialog (4-6 Zeilen).\\
Person A sucht etwas, Person B hilft.
\end{block}

\vspace{0.5cm}

\begin{block}{Beispiel}
A: \es{¿Hay un bolígrafo?}\\
B: \es{Sí, hay uno.}\\
A: \estar{¿Dónde está?}\\
B: \estar{Está encima de la mesa.}
\end{block}

\end{frame}

% ---------------------------------------------------------------
% ZUSAMMENFASSUNG
% ---------------------------------------------------------------
\begin{frame}{Zusammenfassung}

\begin{columns}[T]
\begin{column}{0.48\textwidth}
\begin{block}{\es{hay} = es gibt}
\begin{itemize}
    \item + un/una
    \item + Zahl
    \item + mucho/poco
    \item + ohne Artikel
\end{itemize}
\de{Existenz (unbestimmt)}
\end{block}
\end{column}

\begin{column}{0.48\textwidth}
\begin{block}{\estar{estar} = sein (Ort/Zustand)}
\begin{itemize}
    \item + el/la (bestimmt)
    \item + Präposition
    \item estoy, estás, está...
\end{itemize}
\de{Ort (bestimmt)}
\end{block}
\end{column}
\end{columns}

\vspace{0.5cm}

\begin{center}
{\Large \es{hay} $\neq$ \estar{estar}}\\[0.3em]
{\small Existenz $\neq$ Ort}
\end{center}

\end{frame}

% ---------------------------------------------------------------
% SCHLUSSFOLIE
% ---------------------------------------------------------------
{
\setbeamercolor{background canvas}{bg=spanischrot}
\begin{frame}[plain]
\begin{center}
\vspace{3cm}
{\Huge\bfseries\textcolor{white}{¡Muy bien!}}\\[1cm]
{\Large\textcolor{white}{Bis zur nächsten Stunde!}}
\end{center}
\end{frame}
}

\end{document}
